\section{Explicit level-2 conformal blocks}
\label{sec:blocks}

Let $W=\Inv((\C^2)^{\otimes4})$ have the orthonormal basis
\begin{equation}
 e_0=s_{01}s_{23},\qquad
 e_1=\frac{t^+_{01}t^-_{23}+t^-_{01}t^+_{23}-t^0_{01}t^0_{23}}{\sqrt3},
\end{equation}
where $s=(|\uparrow\downarrow\rangle-|\downarrow\uparrow\rangle)/\sqrt2$
and $t^+,t^0,t^-$ are the normalized triplets. Use
$t^a=\sigma^a/\sqrt2$ and $T_{ij}=\sum_a t_i^at_j^a$. Put
$\nu=(k+2)^{-1}=1/4$. Then $C_{1/2}=3/2$ and
$h=\nu C_{1/2}/2=3/16$. In this basis,
\begin{equation}
 D=T_{01}=\begin{pmatrix}-3/2&0\\0&1/2\end{pmatrix},\quad
 C=T_{02}=\begin{pmatrix}0&-\sqrt3/2\\-\sqrt3/2&-1\end{pmatrix},\quad
 B=T_{12}=\begin{pmatrix}0&\sqrt3/2\\\sqrt3/2&-1\end{pmatrix}.
 \label{eq:matrices}
\end{equation}
In particular $D+C+B=-3I/2$. Replacing these generators by $\sigma^a/2$
without changing the KZ coefficient would incorrectly halve the vacuum
OPE exponent.

Translation and scaling covariance give
\begin{equation}
 F(0,a,b,\infty)=b^{-2h}f(r),\qquad r=a/b,\qquad
 f'(r)=\nu\left(\frac Dr+\frac B{r-1}\right)f(r).
 \label{eq:kz}
\end{equation}
The first adjacent pair in \eqref{eq:labels} has label $0$ on both outside
intervals. Its OPE lies in $\mathcal H_{00}=V_0$ and therefore selects
\begin{equation}
 f^{(0)}(r)=r^{-3/8}(e_0+O(r)),
 \label{eq:selection}
\end{equation}
with no independent spin-one Frobenius branch. This selection is fixed by
the boundary labels, rather than by choosing initial data for favorable
signs. A single selected block generally has both tensor components.

For $0<r<1$, write $s=\sqrt{1-r}$ and $q=[r(1-r)]^{-3/8}$, using positive
real roots. The two unit-leading Frobenius solutions are
\begin{align}
 f^{(0)}(r)&=\frac{q}{2\sqrt2}
 \begin{pmatrix}(1+s)^{3/2}\\-\sqrt3(1-s)\sqrt{1+s}\end{pmatrix},
 \label{eq:block0}\\
 f^{(1)}(r)&=\frac{q}{\sqrt6}
 \begin{pmatrix}-(1-s)^{3/2}\\\sqrt3(1+s)\sqrt{1-s}\end{pmatrix}.
 \label{eq:block1}
\end{align}
The second solution has leading behavior $r^{1/8}e_1$. Both columns are
useful for transport and fusion, although the specified boundary conditions
select \eqref{eq:block0}.

For a direct derivation, eliminate the second component of \eqref{eq:kz}
and write the first as $qP(r)$. The scalar equation becomes
\begin{equation}
 r(1-r)P''-\tfrac12P'-\tfrac3{16}P=0.
\end{equation}
After $s=\sqrt{1-r}$ it is
$(1-s^2)P_{ss}+sP_s-3P/4=0$, solved by $(1+s)^{3/2}$ and
$(1-s)^{3/2}$. Substitution in the first KZ component and normalization of
the leading tensors give \eqref{eq:block0}--\eqref{eq:block1}.

An independent local construction writes
$f=r^\lambda\sum_{n\ge0}v_nr^n$ and uses
\begin{equation}
 [(\lambda+n)I-\nu D]v_n
 =-\nu B\sum_{m=0}^{n-1}v_m,\qquad n\ge1.
 \label{eq:recurrence}
\end{equation}
Take $(\lambda,v_0)=(-3/8,e_0)$ or $(1/8,e_1)$. The exponent difference
is $1/2$, so the displayed inverses exist. For example,
\begin{equation}
 f^{(0)}(r)=r^{-3/8}\left[
 \binom10+r\binom0{-\sqrt3/4}
 +r^2\binom{3/64}{-\sqrt3/8}+O(r^3)\right].
\end{equation}
At $(u,x,y)=(0,1,2)$, the physical selected state is
\begin{equation}
 M_0=2^{-3/8}f^{(0)}(1/2)
 \simeq (1.02266239412001,-0.30390758736355)^T.
 \label{eq:selected_initial}
\end{equation}
